\section{Conclusion}
\label{sec:conclusion}

Every agent-memory defence is itself an attack surface, because every
defence decides on a signal an adversary can reach. We defined the
curation-targeted attack, in which no poison is injected and the curator
is instead fed corrupted evidence until it starves the memory it was
protecting, and measured six mechanisms under it. They separate on one
property: whether a false signal costs a life or costs energy that can be
won back. A bounded credit buffer with earn-back survives, and then fails
too, between two and four lies per cycle. We publish that boundary
because a defence whose limit is unstated is a defence nobody can deploy
safely.

The ledger we defend with, \textsc{darwin-memo}, needs no reward model,
no judge in its core update: settlement consumes an outcome report.
The reported synthetic update timings exclude the cost of producing that
report and do not establish end-to-end operating savings. We are precise about how far that carries. Outside an
adversarial regime a one-line heuristic ties it, a confidence bandit
matches it under one-sided noise, and on real software-engineering tasks
it shows no advantage over a token-matched random control. The regime map
is our second contribution precisely because those losses are real: a
reader deciding what to deploy is better served by knowing which axis
they are on than by a claim of dominance our own tables would refute.

What we would most like tested next is the threat model rather than the
mechanism: if curation-targeted attacks are as general as they appear,
every memory defence owes an answer to what happens when its deciding
signal is the thing under attack. All code, data and protocols are
released with manifest-checked reproducibility.
